%!TEX root = ../main.tex
% =============================================================================
%  摘要专用页（格式规范第三条：含标题与关键词，无需英文，原则上不超过一页；
%              论文自本页开始编页码，页脚中部，阿拉伯数字从 1 起）
%  本页全部数字取自 intermediate/paper_numbers.json 与 v2 新增对照实验
%  写法：背景一句 → 逐问"结构理由 + 模型 + 量化结果" → 验证一句 → 结论与边界一句
% =============================================================================
{\linespread{1.35}\selectfont
\begin{center}
  {\zihao{3}\heiti \papertitle\par}
\end{center}

\begin{center}
  {\zihao{4}\heiti 摘\quad 要}
\end{center}

并网微网在分时电价与光伏出力双重不确定性下运行：光伏盈余落在电价低谷的午间，负载高峰落在电价尖峰的
傍晚，而储能容量与充放功率均有限，收益只能来自把电量从低价时段搬到高价时段。本文按信息结构由弱到强给
出四层输出，每层只增加一个结构性对象，并用三组对照实验检验其结论。

\textbf{问题 1} 中电价、负载与光伏预测给定且要求两端储电量相等，目标与约束均为线性，故建立
\textbf{线性规划}模型；由\textbf{对偶变量}给出储电水价值，在储电量界不活跃时由互补松弛得到最优充放电的
\textbf{阈值型刻画}（死区比 \(1/(\eta_c\eta_d)=1.2346\)）。最优全天购电费 \textbf{35\,126.95 元}，
相对无储能基线 \(48\,052.05\) 元节省 \(26.90\%\)，充、放电量 \(20\,740.67\) 与 \(16\,799.94\)~kWh
（比值 \(1.2346\) 与往返效率互为倒数），弃光为 0；容量摆幅是主要瓶颈（解除容量上限可再省
\(22.14\%\)，解除功率上限仅 \(1.66\%\)）。

\textbf{问题 2} 的计划依据为 \(0{:}00\) 光伏预报、结算依据为当日实际值，缺口按 \(5\) 倍电价紧急购电，
故建立\textbf{两阶段随机规划}模型（场景集按日总残差分层取 medoid、一次存盘固定），并用
\textbf{报童型临界分位数}给出安全边际的解析近似（缺货概率两界 \(1/10\)--\(1/5\)）。334 天总费用
\textbf{14\,264\,792.13 元}，紧急购电量 \(172\,015.82\)~kWh、出现在 296 天，节省 \(13.06\%\)；
安全边际的\textbf{边际价值 3\,133\,767.55 元}（消融：移除后费用升至 \(17\,398\,559.68\) 元、
紧急购电量增至 \(6.72\) 倍）。

\textbf{问题 3} 有四个预报时刻，故建立含\textbf{非预期性}的\textbf{多阶段滚动优化}模型（分段线性凸
偏差费用被精确线性化为 LP，无需 0--1 变量）。费用降至 \textbf{13\,932\,402.95 元}，信息价值
\(332\,389.18\) 元（\(2.33\%\)），已实现完全信息上界价值的 \(85.1\%\)；\textbf{再加密最多再省
\(0.42\%\)}，故不建议继续加密预报时刻，应把投入用于提高现有四次预报（尤其上午）的精度。

\textbf{问题 4} 中电价本身随机：先用\textbf{C 藤 pair-Copula}刻画相依结构（电价—负载
\(\tau=+0.698\)；电价—光伏由自由度 \(3.23\) 的 \(t\) Copula 给出 \(\tau=-0.238\)、对称尾相依
\(0.0479\)，Gaussian 族会把它抹平），再用\textbf{分时骨架加随机残差}两层结构建模电价，用保留联合
结构的\textbf{联合情景}与\textbf{条件风险价值}重算问题 2、3：总费用 \textbf{15\,094\,255.84 元}与
\textbf{14\,722\,153.17 元}，节省 \(11.35\%\) 与 \(13.54\%\)；电价波动使费用上升 \(5.81\%/5.67\%\)
且不能由价格水平解释，说明\textbf{高价与高净负荷的共动}才是成本来源。

\textbf{三组对照给出否证性结论。} ① 四种情景生成器在样本外 184 天上的总费用极差仅 \(0.37\%\)，故主
结果对生成范式稳健，但紧急购电量相差 \(48.5\%\)，引用该分项必须带情景口径。② 机会约束最优档
（\(\beta=0.10\)）仍比随机规划贵 \(0.52\%\)，且实测缺电频率约为标称 \(\beta\) 之半，与报童两界互证。
③ 盒式分布鲁棒最优档在\textbf{全部五项指标}上更差（总费用高 \(11.43\%\)、\(\mathrm{CVaR}\) 高
\(8.72\%\)），其价值只在于给出最坏情形上界。分时段风险画像进一步显示风险是结构性的：费用尾部集中在
\(15\)--\(18\) 时，缺电频率峰值出现在 \(11\) 时（\(23.8\%\)）与 \(5\) 时。

\textbf{验证。} 全部结果通过原始可行性（供给缺口残差 \(3.41\times10^{-13}\)~kWh）、最优性（阈值律
一致性，144 槽违反次数为 0）与费用独立复算（相对误差 \(\le10^{-6}\)）三重检验；审计 A1--A10 在
10~min 级解上全部通过，越界、互斥与负购电计数均为 0；另用\textbf{L 形（Benders）分解}这条独立数值
路径复核最优值（334 天最大相对间隙 \(3.76\times10^{-15}\)）。

\textbf{结论与边界。} 在本题参数范围内储能的核心价值是\textbf{把午间光伏盈余与低谷电量搬运到傍晚高价
时段}；当峰谷价比高于盈亏平衡阈值 \(1/(\eta_c\eta_d)\approx1.235\) 时套利有利。结论适用于价格接受型、
不外送、单一储能的并网微网，且节省幅度按电能量购电费用口径计算，是\textbf{保守下界}；允许上网售电、
引入需求响应或计入容量与辅助服务收益时须重新标定。

\noindent{\heiti 关键词：}\quad 微网；储能调度；线性规划；两阶段随机规划；报童模型；多阶段滚动优化；
Copula 联合情景；CVaR；信息价值
}
\newpage
